\begin{textsample}{POS Dim 1 – 2020 – Score 21.00 – 2020/t004235.txt}  \label{ex:f1_pos_073}
Chris Anderson : Dr.
Jane Goodall, welcome.
Jane Goodall : Thank you, and I think, you know, we couldn ' t have a complete interview unless people know Mr.
H is with me, because everybody knows Mr.
H.
CA : Hello, Mr.
H.
In your TED Talk 17 years ago, you warned us about \textbf{the} dangers of humans crowding out \textbf{the} natural world.
Is there any sense in which you feel that \textbf{the} current pandemic is kind of, nature striking back?
JG : It ' s very, very clear that these zoonotic diseases, like \textbf{the} corona and HIV / AIDS and all sorts of other diseases that we catch from animals, that ' s partly to do with destruction of \textbf{the} environment, which, as animals lose \textbf{habitat}, they get crowded together and sometimes that means that a virus from a reservoir species, where it ' s lived harmoniously for maybe hundreds of years, jumps into a new species, then you also get animals being pushed into closer contact with humans.
And sometimes one of these animals that has caught a virus can -- you know, provides \textbf{the} opportunity for that virus to jump into people and create a new disease, like COVID-19.
And in addition to that, we are so disrespecting animals.
We hunt them, we kill them, we eat them, we traffic them, we send them off to \textbf{the} wild-animal markets in Asia, where they ' re in terrible, cramped conditions, in tiny cages, with people being contaminated with \textbf{blood} and urine and feces, ideal conditions for a virus to \textbf{spill} from an animal to an animal, or an animal to a person.
CA : I ' d love to just dip \textbf{backwards} in time for a bit, because your story is so extraordinary.
I mean, despite \textbf{the} arguably even more sexist attitudes of \textbf{the} 1960s, somehow you were able to break through and become one of \textbf{the} world ' s leading scientists, discovering this astonishing series of facts about chimpanzees, such as their tool use and so much more.
What was it about you, do you think, that allowed you to make such a breakthrough?
JG : Well, \textbf{the} thing is, I was born loving animals, and \textbf{the} most important thing was, I had a very supportive mother.
She didn ' t get mad when she found earthworms in my bed, she just said they better be in \textbf{the} garden.
And she didn ' t get mad when I \textbf{disappeared} for four hours and she called \textbf{the} police, and I was sitting in a hen house, because nobody would tell me where \textbf{the} hole was where \textbf{the} egg came out.
I had no dream of being a scientist, because women didn ' t do that sort of thing.
In fact, there weren ' t any man doing it back then, either.
And everybody laughed at me except Mom, who said,"If you really want this, you ' re going to have to work awfully hard, take advantage of every opportunity, if you don ' t give up, maybe you ' ll find a way."CA : And somehow, you were able to kind of, earn \textbf{the} trust of chimpanzees in \textbf{the} way that no one else had.
Looking back, what were \textbf{the} most exciting moments that you discovered or what is it that people still don ' t get about chimpanzees?
JG : Well, \textbf{the} thing is, you say,"See things nobody else had, get their trust."Nobody else had tried.
Quite honestly.
So, basically, I used \textbf{the} same \textbf{techniques} that I had to study \textbf{the} animals around my home when I was a child.
Just sitting, patiently, not trying to get too close too quickly, but it was awful, because \textbf{the} money was only for six months.
I mean, you can imagine how difficult to get money for a young girl with no degree, to go and do something as bizarre as sitting in a forest.
And you know, finally, we got money for six months from an American philanthropist, and I knew with time I ' d get \textbf{the} chimps ' trust, but did I have time?
And weeks became months and then finally, after about four months, one chimpanzee began to lose his fear, and it was he that on one occasion I saw -- I still wasn ' t really close, but I had my binoculars -- and I saw him using and making tools to \textbf{fish} for termites.
And although I wasn ' t terribly surprised, because I ' ve read about things captive chimps could do -- but I knew that science believed that humans, and only humans, used and made tools.
And I knew how excited [ Dr.
Louis ] Leakey would be.
And it was that observation that enabled him to go to \textbf{the} National Geographic, and they said,"OK, we ' ll continue to support \textbf{the} research,"and they sent Hugo van Lawick, \textbf{the} photographer-filmmaker, to record what I was seeing.
So a lot of scientists didn ' t want to believe \textbf{the} tool-using.
In fact, one of them said I must have taught \textbf{the} chimps.
Since I couldn ' t get near them, it would have been a miracle.
But anyway, once they saw Hugo ' s film and that with all my descriptions of their behavior, \textbf{the} scientists had to start changing their minds.
CA : And since then, numerous other discoveries that placed chimpanzees much closer to humans than people cared to believe.
I think I saw you say at one point that they have a sense of humor.
How have you seen that expressed?
JG : Well, you see it when they ' re playing games, and there ' s a bigger one playing with a little one, and he ' s trailing a vine around a tree.
And every time \textbf{the} little one is about to catch it, \textbf{the} bigger one pulls it away, and \textbf{the} little one starts crying and \textbf{the} big one starts laughing.
So, you know.
CA : And then, Jane, you observed something much more troubling, which was these instances of chimpanzee gangs, tribes, groups, being brutally violent to each other.
I ' m curious how you process that.
And whether it made you, kind of, I don ' t know, depressed about us, we ' re close to them, did it make you feel that violence is irredeemably part of all \textbf{the} great apes, somehow?
JG : Well, it obviously is.
And my first encounter with human, what I call evil, was \textbf{the} end of \textbf{the} war and \textbf{the} pictures from \textbf{the} Holocaust.
And you know, that really shocked me.
That changed who I was.
I was 10, I think, at \textbf{the} time.
And when \textbf{the} chimpanzees, when I realized they have this dark, brutal side, I thought they were like us but nicer.
And then I realized they ' re even more like us than I had thought.
And at that time, in \textbf{the} early ' 70s, it was very strange, aggression, there was a big thing about, is aggression innate or learned.
And it became political.
And it was, I don ' t know, it was a very strange time, and I was coming out, saying,"No, I think aggression is definitely part of our inherited repertoire of behaviors."And I asked a very respected scientist what he really thought, because he was coming out on \textbf{the} clean slate, aggression is learned, and he said,"Jane, I ' d rather not talk about what I really think."That was a big shock as far as science was concerned for me.
CA : I was brought up to believe a world of all things bright and beautiful.
You know, numerous beautiful films of butterflies and bees and flowers, and you know, nature as this gorgeous landscape.
And many environmentalists often seem to take \textbf{the} stance,"Yes, nature is pure, nature is beautiful, humans are bad,"but then you have \textbf{the} kind of observations that you see, when you actually look at any part of nature in more detail, you see things to be terrified by, honestly.
What do you make of nature, how do you think of it, how should we think of it?
JG : Nature is, you know, I mean, you think of \textbf{the} whole spectrum of \textbf{evolution}, and there ' s something about going to a pristine place, and Africa was very pristine when I was young.
And there were animals everywhere.
And I never liked \textbf{the} fact that \textbf{lions} killed, they have to, I mean, that ' s what they do, if they didn ' t kill animals, they would die.
And \textbf{the} big difference between them and us, I think, is that they do what they do because that ' s what they have to do.
And we can plan to do things.
Our plans are very different.
We can plan to cut down a whole forest, because we want to sell \textbf{the} timber, or because we want to build another shopping mall, something like that.
So our destruction of nature and our warfare, we ' re capable of evil because we can sit comfortably and plan \textbf{the} torture of somebody far away.
That ' s evil.
Chimpanzees have a sort of primitive war, and they can be very aggressive, but it ' s of \textbf{the} moment.
It ' s how they feel.
It ' s response to an emotion.
CA : So your observation of \textbf{the} sophistication of chimpanzees doesn ' t go as far as what some people would want to say is \textbf{the} sort of \textbf{the} human superpower, of being able to really simulate \textbf{the} future in our minds in great detail and make long-term plans.
And act to encourage each other to achieve those long-term plans.
That that feels, even to someone who spent so much time with chimpanzees, that feels like a fundamentally different skill set that we just have to take responsibility for and use much more wisely than we do.
JG : Yes, and I personally think, I mean, there ' s a lot of discussion about this, but I think it ' s a fact that we developed \textbf{the} way of communication that you and I are using.
And because we have words, I mean, animal communication is way more sophisticated than we used to think.
And chimpanzees, gorillas, orangutans can learn human sign language of \textbf{the} Deaf.
But we sort of grow up speaking whatever language it is.
So I can tell you about things that you ' ve never heard of.
And a chimpanzee couldn ' t do that.
And we can teach our children about \textbf{abstract} things.
And chimpanzees couldn ' t do that.
So yes, chimpanzees can do all sorts of clever things, and so can \textbf{elephants} and so can crows and so can octopuses, but we design rockets that go off to another \textbf{planet} and little robots taking photographs, and we ' ve designed this extraordinary way of you and me talking in our different parts of \textbf{the} world.
When I was young, when I grew up, there was no TV, there were no cell phones, there was no \textbf{computers}.
It was such a different world, I had a pencil, pen and notebook, that was it.
CA : So just going back to this question about nature, because I think about this a lot, and I struggle with this, honestly.
So much of your work, so much of so many people who I respect, is about this passion for trying not to screw up \textbf{the} natural world.
So is it possible, is it healthy, is it essential, perhaps, to simultaneously accept that many aspects of nature are \textbf{terrifying}, but also, I don ' t know, that it ' s awesome, and that some of \textbf{the} awesomeness comes from its potential to be \textbf{terrifying} and that it is also just breathtakingly beautiful, and that we cannot be ourselves, because we are part of nature, we cannot be whole unless we somehow embrace it and are part of it?
Help me with \textbf{the} language, Jane, on how that relationship should be.
JG : Well, I think one of \textbf{the} problems is, you know, as we developed our intellect, and we became better and better at modifying \textbf{the} environment for our own use, and creating fields and growing crops where it used to be forest or woodland, and you know, we won ' t go into that now, but we have this ability to change nature.
And as we ' ve moved more into towns and cities, and relied more on technology, many people feel so divorced from \textbf{the} natural world.
And there ' s hundreds, thousands of children growing up in inner cities, where there basically isn ' t any nature, which is why this movement now to green our cities is so important.
And you know, they ' ve done experiments, I think it was in Chicago, I ' m not quite sure, and there were various empty lots in a very violent part of town.
So in some of those areas they made it green, they put trees and flowers and things, shrubs in these vacant lots.
And \textbf{the} crime rate went right down.
So then of course, they put trees in \textbf{the} other half.
So it just shows, and also, there have been studies done showing that children really need green nature for good psychological development.
But we are, as you say, part of nature and we disrespect it, as we are, and that is so terrible for our children and our children ' s children, because we rely on nature for clean air, clean water, for regulating climate and rainfall.
Look what we ' ve done, look at \textbf{the} climate crisis.
That ' s us.
We did that.
CA : So a little over 30 years ago, you made this shift from scientist mainly to activist mainly, I guess.
Why?
JG : Conference in 1986, scientific one, I ' d got my PhD by then and it was to find out how chimp behavior differed, if it did, from one environment to another.
There were six study sites across Africa.
So we thought, let ' s bring these scientists together and explore this, which was fascinating.
But we also had a session on conservation and a session on conditions in some captive situations like medical research.
And those two sessions were so shocking to me.
I went to \textbf{the} conference a a scientist, and I left as an activist.
I didn ' t make \textbf{the} decision, something happened inside me.
CA : So you spent \textbf{the} last 34 years sort of tirelessly campaigning for a better relationship between people and nature.
What should that relationship look like?
JG : Well, you know, again you come up with all these problems.
People have to have space to live.
But I think \textbf{the} problem is that we ' ve become, in \textbf{the} affluent societies, too greedy.
I mean, honestly, who needs four houses with huge grounds?
And why do we need yet another shopping mall?
And so on and so on.
So we are looking at short-term economic benefit, money has become a sort of god to worship, as we lose all spiritual connection with \textbf{the} natural world.
And so we ' re looking for short-term monetary gain, or power, rather than \textbf{the} health of \textbf{the} \textbf{planet} and \textbf{the} future of our children.
We don ' t seem to care about that anymore.
That ' s why I ' ll never stop fighting.
CA : I mean, in your work specifically on chimpanzee conservation, you ' ve made it practice to put people at \textbf{the} center of that, local people, to engage them.
How has that worked and do you think that ' s an essential idea if we ' re to succeed in protecting \textbf{the} \textbf{planet}?
JG : You know, after that \textbf{famous} conference, I thought, well, I must learn more about why chimps are vanishing in Africa and what ' s happening to \textbf{the} forest.
So I got a bit of money together and went out to visit six range countries.
And learned a lot about \textbf{the} problems faced by chimps, you know, hunting for bushmeat and \textbf{the} live animal trade and caught in snares and human populations growing and needing more land for their crops and their cattle and their villages.
But I was also learning about \textbf{the} plight faced by so many people. \textbf{The} absolute poverty, \textbf{the} lack of health and education, \textbf{the} degradation of \textbf{the} land.
And it came to a head when I flew over \textbf{the} tiny Gombe National Park.
It had been part of this equatorial forest belt right across Africa to \textbf{the} west coast, and in 1990, it was just this little \textbf{island} of forest, just tiny national park.
All around, \textbf{the} hills were bare.
And that ' s when it hit me.
If we don ' t do something to help \textbf{the} people find ways of living without destroying their environment, we can ' t even try to save \textbf{the} chimps.
So \textbf{the} Jane Goodall Institute began this program"Take Care,"we call it"TACARE."And it ' s our method of community-based conservation, totally holistic.
And we ' ve now put \textbf{the} tools of conservation into \textbf{the} hand of \textbf{the} villagers, because most Tanzanian wild chimps are not in protected areas, they ' re just in \textbf{the} village forest reserves.
And so, they now go and measure \textbf{the} health of their forest.
They ' ve understood now that protecting \textbf{the} forest isn ' t just for wildlife, it ' s their own future.
That they need \textbf{the} forest.
And they ' re very proud. \textbf{The} volunteers go to workshops, they learn how to use \textbf{smartphones}, they learn how to upload into platform and \textbf{the} cloud.
And so it ' s all transparent.
And \textbf{the} trees have come back, there ' s no bare hills anymore.
They agreed to make a buffer zone around Gombe, so \textbf{the} chimps have more forest than they did in 1990.
They ' re opening up corridors of forest to link \textbf{the} scattered chimp groups so that you don ' t get too much inbreeding.
So yes, it ' s worked, and it ' s in six other countries now.
Same thing.
CA : I mean, you ' ve been this extraordinary tireless voice, all around \textbf{the} world, just traveling so much, speaking everywhere, inspiring people everywhere.
How on earth do you find \textbf{the} energy, you know, \textbf{the} fire to do that, because that is exhausting to do, every meeting with lots of people, it is just physically exhausting, and yet, here you are, still doing it.
How are you doing this, Jane?
JG : Well, I suppose, you know, I ' m obstinate, I don ' t like giving up, but I ' m not going to let these CEOs of big companies who are destroying \textbf{the} forests, or \textbf{the} politicians who are unraveling all \textbf{the} protections that were put in place by previous presidents, and you know who I ' m talking about.
And you know, I ' ll go on fighting, I care about, I ' m passionate about \textbf{the} wildlife.
I ' m passionate about \textbf{the} natural world.
I love forests, it hurts me to see them damaged.
And I care passionately about children.
And we ' re stealing their future.
And I ' m not going to give up.
So I guess I ' m blessed with good genes, that ' s a gift, and \textbf{the} other gift, which I discovered I had, was communication, whether it ' s writing or speaking.
And so, you know, if going around like this wasn ' t working, but every time I do a \textbf{lecture}, people come up and say,"Well, I had given up, but you ' ve inspired me, I promise to do my bit."And we have our youth program"Roots and Shoots"now in 65 countries and growing fast, all ages, all choosing projects to help people, animals, \textbf{the} environment, rolling up their sleeves and taking action.
And you know, they look at you with shining eyes, wanting to tell Dr.
Jane what they ' ve been doing to make \textbf{the} world a better place.
How can I let them down?
CA : I mean, as you look at \textbf{the} \textbf{planet} ' s future, what worries you most, actually, what scares you most about where we ' re at?
JG : Well, \textbf{the} fact that we have a small window of time, I believe, when we can at least start healing some of \textbf{the} harm and slowing down climate change.
But it is closing, and we ' ve seen what happens with \textbf{the} lockdown around \textbf{the} world because of COVID-19 : clear skies over cities, some people breathing clean air that they ' ve never breathed before and looking up at \textbf{the} shining skies at night, which they ' ve never seen properly before.
And you know, so what worries me most is how to get enough people, people understand, but they ' re not taking action, how to get enough people to take action?
CA : National Geographic just launched this extraordinary film about you, highlighting your work over six decades.
It ' s titled"Jane Goodall : \textbf{The} Hope."So what is \textbf{the} hope, Jane?
JG : Well, \textbf{the} hope, my greatest hope is all these young people.
I mean, in China, people will come up and say,"Well, of course I care about \textbf{the} environment, I was in ' Roots and Shoots ' in primary school."And you know, we have"Roots and Shoots"just hanging on to \textbf{the} values and they ' re so enthusiastic once they know \textbf{the} problems and they ' re empowered to take action, they are clearing \textbf{the} streams, removing invasive species humanely.
And they have so many ideas.
And then there ' s, you know, this extraordinary intellect of ours.
We ' re beginning to use it to come up with technology that really will help us to live in greater harmony, and in our individual lives, let ' s think about \textbf{the} consequences of what we do each day.
What do we buy, where did it come from, how was it made?
Did it harm \textbf{the} environment, was it cruel to animals?
Is it cheap because of child slave labor?
Make ethical choices.
Which you can ' t do if you ' re living in poverty, by \textbf{the} way.
And then finally, this indomitable spirit of people who tackle what seems impossible and won ' t give up.
You can ' t give up when you have those...
But you know, there are things that I can ' t fight.
I can ' t fight corruption.
I can ' t fight military regimes and dictators.
So I can only do what I can do, and if we all do \textbf{the} bits that we can do, surely that makes a whole that eventually will win out.
CA : So, last question, Jane.
If there was one idea, one thought, one seed you could plant in \textbf{the} minds of everyone watching this, what would that be?
JG : You know, just remember that every day you live, you make an impact on \textbf{the} \textbf{planet}.
You can ' t help making an impact.
And at least, unless you ' re living in extreme poverty, you have a choice as to what sort of impact you make.
Even in poverty you have a choice, but when we are more affluent, we have a greater choice.
And if we all make ethical choices, then we start moving towards a world that will be not quite so desperate to leave to our great-grandchildren.
That ' s, I think, something for everybody.
Because a lot of people understand what ' s happening, but they feel helpless and hopeless, and what can they do, so they do nothing and they become apathetic.
And that is a huge danger, apathy.
CA : Dr.
Jane Goodall, wow.
I really want to thank you for your extraordinary life, for all that you ' ve done and for spending this time with us now.
Thank you.
JG : Thank you.

% matched lemmas: abstract, backwards, blood, computer, disappear, elephant, evolution, famous, fish, habitat, island, lecture, lion, planet, smartphone, spill, technique, terrifying, the
\end{textsample}
